\chapter{Componenti Base}
\label{chap:componenti_base}

In questo capitolo vengono descritti i blocchi aritmetici fondamentali progettati ed implementati per costituire la base delle elaborazioni successive.

\section{Full Adder}
L'elemento atomico per le operazioni aritmetiche è il Full Adder, un circuito combinatorio in grado di calcolare la somma di tre bit: due addendi ($A$, $B$) e un riporto in ingresso ($C_{in}$).



\subsection{Logica Implementativa}
Per l'implementazione su FPGA è stata scelta una descrizione \textit{behavioral}. Sebbene le equazioni canoniche siano:
\begin{equation}
	S = A \oplus B \oplus C_{in}
\end{equation}
\begin{equation}
	C_{out} = (A \cdot B) + (C_{in} \cdot (A \oplus B))
\end{equation}

Il codice VHDL sfrutta il segnale di propagazione $P$ per ottimizzare l'inferenza della logica di riporto (Carry Logic), fondamentale per le prestazioni temporali:
\begin{itemize}
	\item Se $P = 1$ ($A \neq B$), il riporto $C_{out}$ è determinato da $C_{in}$.
	\item Se $P = 0$ ($A = B$), il riporto $C_{out}$ è uguale ad $A$.
\end{itemize}

\begin{lstlisting}[language=VHDL, caption={Entity e Architecture del Full Adder}, label={lst:full_adder}]
	entity full_adder is
		port (
			a, b, cin : in  STD_LOGIC;
			s, cout   : out STD_LOGIC
		);
	end entity full_adder;
	
	architecture behavioral of full_adder is
		signal p : STD_LOGIC;
		begin
		-- Calcolo Propagate
		p <= a xor b;
		s <= p xor cin;
	
		-- Mux-based Carry logic (Efficiente su FPGA)
		cout <= cin when p = '1' else a;
	end architecture behavioral;
\end{lstlisting}

\section{N-Bit Ripple Carry Adder (RCA)}
Per sommare vettori di bit, la struttura più semplice è il Ripple Carry Adder, ottenuto collegando in cascata $N$ istanze di Full Adder.



\subsection{Architettura Parametrica}
Il componente è stato reso generico attraverso il parametro `N`. L'uso del costrutto \texttt{generate} permette di istanziare dinamicamente la catena di addizionatori in base alla larghezza di bit desiderata a tempo di sintesi.

\begin{itemize}
	\item \textbf{Vantaggi:} Semplicità implementativa e occupazione d'area ridotta.
	\item \textbf{Svantaggi:} Ritardo di propagazione lineare ($T_{delay} \propto N$), poiché il riporto deve attraversare l'intera catena dal LSB al MSB.
\end{itemize}

\begin{lstlisting}[language=VHDL, caption={Ripple Carry Adder parametrico}, label={lst:rca}]
	entity ripple_carry_adder is
		generic ( N : POSITIVE ); 
		port ( 
			a, b : in  std_logic_vector(N-1 downto 0);
			cin  : in  std_logic;
			s    : out std_logic_vector(N-1 downto 0);
			cout : out std_logic
		);
	end entity ripple_carry_adder;
	
	architecture structural of ripple_carry_adder is
		signal c : std_logic_vector(N downto 0);
		begin
		c(0) <= cin; -- Carry in iniziale
	
		loop_n: for i in 0 to N - 1 generate
			fa_i: entity work.full_adder
			port map (
				a    => a(i),
				b    => b(i),
				cin  => c(i),
				s    => s(i),
				cout => c(i + 1) -- Chain connection
			);
		end generate loop_n;
	
		cout <= c(N); -- Ultimo riporto
	end architecture structural;
\end{lstlisting}

\section{Strategie di Testing}
La verifica dei componenti aritmetici ha seguito due approcci distinti in base alla complessità dello spazio degli stati.

\subsection{Unit Testing: Full Adder (Exhaustive)}
Poiché il Full Adder ha solo 3 ingressi ($2^3 = 8$ casi totali), è stato possibile implementare un test esaustivo iterando su una \textbf{Truth Table} completa. La testbench è \textit{self-checking}: confronta l'uscita del DUT (Device Under Test) con il valore atteso memorizzato nella tabella.

\begin{lstlisting}[language=VHDL, caption={Testbench Full Adder con Tabella della Verità}, label={lst:tb_fa}]
	type truth_table_type is array (0 to 7) of std_logic_vector(4 downto 0);
	
	constant TRUTH_TABLE : truth_table_type := (
	"000" & "00", -- 0+0+0 = 0 (C=0)
	"011" & "10", -- 0+1+1 = 0 (C=1)
	"111" & "11", -- 1+1+1 = 1 (C=1)
	-- ... altri casi ...
	others => (others => '0')
	);
	
	-- Loop di simulazione
	-- (Codice omesso per brevita', itera su TRUTH_TABLE)
\end{lstlisting}

\subsection{Unit Testing: Ripple Carry Adder (Directed)}
Per l'addizionatore a $N$ bit, verificare tutte le combinazioni ($2^{2N+1}$) è impraticabile. Si è optato per un \textit{Directed Testing} focalizzato sui "Corner Cases" (casi limite) critici per la propagazione del riporto.

\begin{lstlisting}[language=VHDL, caption={Corner Cases per RCA}, label={lst:tb_rca_cases}]
	constant TRUTH_TABLE : truth_table_type := (
	-- Caso Base: 5 + 5 = 10
	2 => std_logic_vector(x"05") & std_logic_vector(x"05") & '0' & ...
	
	-- Caso Overflow (Carry Out generation):
	-- 255 (FF) + 1 (01) = 0 (00) con Carry=1
	3 => std_logic_vector(x"FF") &
	std_logic_vector(x"01") &
	'0' & -- Cin
	'1' & -- Expected Cout
	std_logic_vector(x"00"), -- Expected Sum
	
	-- Altri casi...
	others => (others => '0')
	);
\end{lstlisting}

Gli scenari validati includono:
\begin{enumerate}
	\item \textbf{Overflow:} Verifica che il riporto esca correttamente dal MSB ($C_{out}=1$).
	\item \textbf{Propagazione Completa:} Verifica il caso peggiore di ritardo, dove un riporto generato al bit 0 deve propagarsi fino all'uscita $C_{out}$ (es. $011..1 + 000..1$).
\end{enumerate}